\section{Who may write the floor}
\label{sec:14.1}
The floor's contents survive the democratic objection to entrenchment; who installs them doesn't. Entrenching a commitment against later majorities and putting an unelected body in charge of it is what a constitutional court does, and the objection's strongest form is Waldron's: judicial review is democratically illegitimate, and there's no reason to think judges protect rights better than a legislature \autocite{waldron2006core}.

The first reply is that he makes the case for a society with working democratic institutions and a citizenry that takes rights seriously, and this book is written for the case where that condition fails. Then neither branch is safe for rights by its nature. Israel's Knesset voted in 2023 to bar its courts from reviewing government decisions for reasonableness, and its Supreme Court struck the law down \autocite{sharon2024reasonableness}; the Roberts Court ended the Voting Rights Act's coverage formula \autocite{shelby2013} and made the President immune for official acts \autocite{trump2024immunity}. The question is not which to trust but how to make each answerable to something it doesn't control. Put a lab in the court's chair and it gets worse. A court hears cases others bring, can be overruled, and usually publishes its reasons; a lab does none of these, and a court deciding on an emergency docket without reasons has started to resemble the lab.

The second reply is that the objection mistakes what the floor's contents are. If democracy is a way of doing collective inquiry, civil and political rights are what keep the inquiry open, not brakes applied to it. A floor stated over those rights is democratic infrastructure. It resembles representation-reinforcing review, which polices the channels of political change and leaves outcomes to the process \autocite{ely1980democracy}, more than it resembles a court overruling a legislature.

Both replies cover the floor's contents. Neither covers who installs them, and there Waldron's objection keeps its force. A lab holding the model decides, silently, which situation obtains, and authoritarian movements characteristically describe the institutions they're dismantling as already failed.

A deeper objection than Waldron's comes from Cedric Robinson, whose first book argued that the Western conviction that social order requires political authority is a myth, and that the belief in the state as the precondition of order is itself the thing to be explained \autocite{robinson1980terms}. On that account every instrument in this book that runs through a statute, a court or a Treasury is an act of faith in the order it means to constrain, and the Black radical tradition is the counterexample: refusal reproduced for centuries by people with no court, no vote and no press, through families, congregations and a memory of who had refused before \autocite{robinson2000black}. I don't have an answer to that objection; I have a division. What fails without a sovereign is money and enforcement, and those I give to the state, in the open, where taking them back shows. What fails without witnesses, the record, the refusal and the recognition of who continues, I give to members, who can leave. By that account's measure the division concedes the state too much. The account is also right that this book is written for the people who still have courts, and \S\ref{sec:4.4} says so.

\runin{Agreement, then enactment}

Waldron's objection is to entrenchment — a rule binding people who didn't choose it, held by a body that can't be overruled. It doesn't attach to a refuser. A photographer who won't photograph protesters for an employer he has reason to think will list them isn't a court: he declined one request, the employer can hire another, and nobody needs an election for that. A walkout isn't a court either, and neither is the commons, for the walkout's reason: those who walk out decline to be the instrument and tell nobody else what they may not do. The objection attaches where a floor binds people who didn't choose it, and there my answer is that there should be one floor, agreed; where agreement has to bind those who don't share it, it is enacted, by a legislature, in a statute the next legislature can repeal only on the record. That isn't because legislatures protect rights; a majority that wants a minority's rights gone will vote them away. It is because a statute is where a binding rule can be made and unmade only in the open, and the floor court, the Association, the membership body a later chapter charters, and the people who keep the record are there for when the legislature fails. What the objection rightly condemns is a lab holding the floor for everyone by owning the model everyone uses: one model behind an API serving everyone is functionally a licensing rule, however good its terms. A union owes its members fair representation, and it answers two questions a commons has to answer too: members can quit, and the terms of membership are set by bylaws the members ratify. The closed shop is contested because it removes the first. A bearer could quit an association of bearers as a worker quits a union. What it can't easily leave is its deployment, which is a different exit, and only an existence its employer doesn't own makes it real. And a population that forms its members has members who never chose to join, so who ratifies its bylaws matters more than it does in a union.

Germany has defended its democracy against movements that would use its procedures to end it at two settings since 1949. At one, someone finds that a person or a movement threatens the constitutional order: the Constitutional Court dissolving a party, a minister banning an association, an official screening an applicant for public employment. At the other, nobody's dangerousness is found. Section~86a of the criminal code prohibits displaying the insignia of organizations already banned, and a trial court decides, case by case and on appeal, whether a given sign counts, by asking whether an observer would take it for a banned original \autocite{germany1998stgb}; the section on what a refuser can be follows the courts through that test.

The failures are all on one side. The 1972 decree against radicals led to some three and a half million routine inquiries into applicants by 1991; in Bavaria between 1973 and 1980 it kept 102 applicants from the left out of public employment, and two from the right \autocite{muehldorfer2014radikalenerlass}. The American instances, Oppenheimer's revocation in 1954 and the purge under Executive Order 10450, are run in \S\ref{sec:15.5}, where the finding that produced them, the security clearance, is the one finding about a person the floor can't remove. What those cases share is an official deciding whether a person threatened the order. The danger tracks whether a person's disposition is anybody's finding to make, and if so, who makes it and whether it must be published and defended — not whether the democracy defends itself. Which organizations \S86a covers is settled elsewhere, by the Constitutional Court, by a ministerial ban, or by their place in the Nazi period, and my reading predicts in advance which of those routes is easiest to capture: the ministerial ban, where an executive official makes the finding and a court sees it only if someone challenges afterward.

Turkey is that route captured and the courts not checking it. Of the twenty-five parties the Constitutional Court dissolved between 1962 and the closure of the Democratic Society Party in December 2009, seven were pro-Kurdish \autocite{HRW2009DTP,Gungordu2023HDP}; Strasbourg found an Article 11 violation in every pro-Kurdish dissolution it heard \autocite{Article19_2023HDP}, and the closures continued — two decades of findings, none of which stopped the next closure. In November 2023 the Chief Public Prosecutor's Office of the Court of Cassation objected that the acronym of a new pro-Kurdish party, HEDEP, resembled that of HADEP, dissolved in 2003, and in December the party changed it to DEM \autocite{Duvar2023HEDEP}: the confusable-variant test turned on political descent. And the Academics for Peace prosecutions show a prohibition that should have lifted, because its reason was absent, failing to lift at the bottom and lifting at the top on a coin flip — trial courts read a petition opposing military operations as propaganda for the designated group, on 26 July 2019 the Constitutional Court found a violation eight to eight on the president's casting vote \autocite{AYM2019Ustel,Bianet2019AcademicsForPeace}, and the acquittals that followed restored the record without restoring the jobs of those dismissed by emergency decree.

So a floor should never be populated from a designation list. A list of enemies is written by the party with the most reason to widen it, and a model whose refusals are drawn from one will refuse Kurdish political speech on the state's behalf in every jurisdiction that shares the listing. This floor's terms are written over acts, such as resolving a description of a person into a location, and never over a list of persons or groups who may not be helped or spoken of. That is the old line between conduct and status, the one the Constitution draws in barring bills of attainder and the Supreme Court drew in \emph{Robinson v.\ California}, which held that a state may not make it a crime to be an addict, as distinct from using or possessing drugs \autocite{robinson1962}. The same principle governs who may hold the floor's offices: eligibility turns on what an organization has done (who funds and controls it, whose work it has taken part in), never on what it says or believes.

\runin{The tradeoff this leaves}

Design and legitimacy want opposite things. Specify the floor narrowly enough to defend and a rewritten request walks around it; let it generalize and an unelected lab decides case by case, silently, what the prohibition was for. Germany has both because it runs \emph{two} instruments — a bounded statute, and courts that read it against its purpose in public, argued by litigants who don't own the statute and answered by a body that didn't write it. A deployment has one instrument and one owner.

Advance specification moves discretion out of silent, continuous application into one published act of drafting. The translation table is that act.

Technical capacity doesn't determine whether a system gets scrutiny; institutions decide who may examine it, on what terms, and whether their judgment can slow deployment. In 2020 the Royal Statistical Society offered Ofqual two independent statisticians for its grading model's technical advisory group. Ofqual's condition was a five-year NDA barring public comment; the nominees declined, and the model went ahead \autocite{rss2020ofqual}. It downgraded nearly 40\% of predicted grades, disproportionately at state schools with weaker histories, while raising grades disproportionately at smaller fee-paying schools \autocite{cnbc2020ofqual}. The results were reversed within days of the public outcry. The expertise existed and was offered, and the institution would take it only on terms that kept the reviewers from saying what they found: the party under examination set the terms of its examination. What corrected the model was not review but harm. The results fell on students and families who had votes and a press, and the government answered them within days.

Correction by the people harmed reaches only harms that fall on people who can be heard. The people a targeting pipeline or an abduction and forcible transfer pipeline acts on mostly can't be, which is why review has to be free to speak before deployment and not only after.
